% ══════════════════════════════════════════════════════════════════════════════
%                         فصل چهارم
%                 سوسیالیسم: ریشه‌ها و شاخه‌ها
% ══════════════════════════════════════════════════════════════════════════════

\chapter{سوسیالیسم: ریشه‌ها و شاخه‌ها}
\label{ch:socialism}

\begin{flushright}
\textit{«تاریخ همه جوامع تاکنون موجود، تاریخ مبارزه طبقاتی است.»}

— کارل مارکس و فردریش انگلس، مانیفست کمونیست (۱۸۴۸)
\end{flushright}

% ══════════════════════════════════════════════════════════════════════════════
%                         مقدمه فصل
% ══════════════════════════════════════════════════════════════════════════════

\lettrine[lines=3, lraise=0.1, nindent=0.5em]{\textcolor{LeftRed}{س}}{وسیالیسم} 
شاید پرنفوذترین ایده سیاسی دو قرن اخیر بوده است. این واژه که نخستین بار در دهه ۱۸۳۰ به کار رفت، چتری شد بر فراز جنبش‌ها و اندیشه‌های بسیار متنوع—از تعاونی‌های مسالمت‌آمیز تا انقلاب‌های خونین، از دولت‌های توتالیتر تا دموکراسی‌های پیشرفته اسکاندیناوی. در این فصل، ریشه‌های سوسیالیسم را می‌کاویم و شاخه‌های متعدد آن را دنبال می‌کنیم.

% ══════════════════════════════════════════════════════════════════════════════
\section{سوسیالیسم چیست؟}
% ══════════════════════════════════════════════════════════════════════════════

\subsection{تعاریف رقیب}

سوسیالیسم از آن مفاهیمی است که تعریف واحدی ندارد. مدافعان و منتقدانش تعاریف متفاوتی از آن ارائه می‌دهند.

\begin{defbox}[تعریف: سوسیالیسم (به طور کلی)]
سوسیالیسم مجموعه‌ای از ایده‌ها و جنبش‌هاست که این ویژگی‌های مشترک را دارند:
\begin{itemize}[nosep]
    \item نقد نابرابری‌های ناشی از سرمایه‌داری
    \item تأکید بر \textbf{همبستگی اجتماعی} در برابر فردگرایی افراطی
    \item خواست \textbf{مالکیت جمعی یا اجتماعی} ابزار تولید (به اشکال مختلف)
    \item هدف: جامعه‌ای برابرتر، انسانی‌تر و دموکراتیک‌تر
\end{itemize}
\end{defbox}

\begin{table}[htbp]
\centering
\caption{تعاریف مختلف سوسیالیسم}
\label{tab:socialism-definitions}
\begin{tabular}{r|r}
\toprule
\textbf{منبع} & \textbf{تعریف} \\
\midrule
مارکس و انگلس & مرحله گذار به کمونیسم؛ دیکتاتوری پرولتاریا \\
\midrule
سوسیال‌دموکرات‌ها & اقتصاد مختلط با دولت رفاه قوی \\
\midrule
آنارشیست‌ها & خودگردانی کارگری بدون دولت \\
\midrule
لیبرتارین‌های راست & هر نوع دخالت دولت در اقتصاد! \\
\midrule
برنی سندرز & آموزش و بهداشت رایگان (سوسیال‌دموکراسی) \\
\bottomrule
\end{tabular}
\end{table}

\subsection{ویژگی‌های مشترک}

با وجود تنوع، می‌توان ویژگی‌هایی یافت که اغلب سوسیالیست‌ها در آن‌ها اشتراک دارند:

\begin{enumerate}[nosep, rightmargin=1.5cm]
    \item \textbf{نقد سرمایه‌داری:} سرمایه‌داری ذاتاً ناعادلانه، استثمارگر یا ناپایدار است
    \item \textbf{تأکید بر طبقه:} جامعه به طبقات تقسیم شده و منافعشان متضاد است
    \item \textbf{همبستگی:} انسان‌ها موجوداتی اجتماعی‌اند، نه صرفاً فردگرا
    \item \textbf{برابری:} هدف نهایی، جامعه‌ای برابرتر است
    \item \textbf{تغییر:} وضع موجود باید تغییر کند (اصلاحی یا انقلابی)
\end{enumerate}

% ══════════════════════════════════════════════════════════════════════════════
\section{سوسیالیسم پیشامارکسی: آرمان‌شهرگرایان}
% ══════════════════════════════════════════════════════════════════════════════

پیش از مارکس، متفکرانی بودند که مارکس و انگلس آن‌ها را «سوسیالیست‌های تخیلی» یا «آرمان‌شهرگرا» (\term{Utopian}) نامیدند. این نام‌گذاری منتقدانه بود، اما امروزه بسیاری از ایده‌های آنان دوباره مورد توجه قرار گرفته‌اند.

\begin{figure}[htbp]
\centering
\begin{tikzpicture}[scale=0.85]

% عنوان
\node[font=\large\bfseries, text=DarkGray] at (0, 7) {\fa{تایم‌لاین: سوسیالیست‌های پیشامارکسی}};

% خط زمان
\draw[line width=4pt, MediumGray] (-7, 5) -- (7, 5);

% رویدادها
% سن‌سیمون
\fill[LeftRedLight] (-5.5, 5) circle (10pt);
\node[above, text=LeftRedLight, font=\small\bfseries] at (-5.5, 5.4) {\fa{۱۷۶۰-۱۸۲۵}};
\node[below, text width=2.5cm, align=center, text=DarkGray, font=\small] at (-5.5, 4.3) {
    \fa{سن‌سیمون}\\
    \fa{تکنوکراسی}
};

% فوریه
\fill[LeftRed] (-2.5, 5) circle (10pt);
\node[above, text=LeftRed, font=\small\bfseries] at (-2.5, 5.4) {\fa{۱۷۷۲-۱۸۳۷}};
\node[below, text width=2.5cm, align=center, text=DarkGray, font=\small] at (-2.5, 4.3) {
    \fa{شارل فوریه}\\
    \fa{فالانستر}
};

% اوون
\fill[LeftRedDark] (0.5, 5) circle (10pt);
\node[above, text=LeftRedDark, font=\small\bfseries] at (0.5, 5.4) {\fa{۱۷۷۱-۱۸۵۸}};
\node[below, text width=2.5cm, align=center, text=DarkGray, font=\small] at (0.5, 4.3) {
    \fa{رابرت اوون}\\
    \fa{تعاونی‌ها}
};

% پرودون
\fill[OrangeAccent] (3.5, 5) circle (10pt);
\node[above, text=OrangeAccent, font=\small\bfseries] at (3.5, 5.4) {\fa{۱۸۰۹-۱۸۶۵}};
\node[below, text width=2.5cm, align=center, text=DarkGray, font=\small] at (3.5, 4.3) {
    \fa{پرودون}\\
    \fa{آنارشیسم}
};

% مانیفست
\fill[LeftRedDark] (6, 5) circle (12pt);
\node[above, text=LeftRedDark, font=\small\bfseries] at (6, 5.5) {\fa{۱۸۴۸}};
\node[below, text width=2.5cm, align=center, text=DarkGray, font=\small] at (6, 4.2) {
    \fa{مانیفست}\\
    \fa{کمونیست}
};

% فلش تقسیم
\draw[line width=2pt, color=LeftRedDark, dashed] (6, 4) -- (6, 2.5);
\node[text=LeftRedDark, font=\small\bfseries] at (6, 2) {\fa{مارکسیسم}};

\end{tikzpicture}
\caption{پیشگامان سوسیالیسم پیش از مارکس}
\label{fig:pre-marx-socialists}
\end{figure}

\subsection{کلود هانری سن‌سیمون}

\begin{personbox}[کلود هانری دو سن‌سیمون (۱۷۶۰-۱۸۲۵)]

\textbf{ملیت:} فرانسوی

\textbf{ایده محوری:} حکومت متخصصان و صنعتگران

\mediumspace

\person{سن‌سیمون}، اشراف‌زاده‌ای که در انقلاب آمریکا جنگیده بود، از اولین نظریه‌پردازان جامعه صنعتی بود. او استدلال کرد:

\begin{itemize}[nosep]
    \item جامعه صنعتی نیاز به مدیریت \textbf{علمی} دارد
    \item «بیکارگان» (اشراف، روحانیون، حقوقدانان) باید کنار بروند
    \item «صنعتگران» (کارگران، مهندسان، کارفرمایان مولد) باید حکومت کنند
    \item هدف: «بهبود وضعیت فقیرترین و پرشمارترین طبقه»
\end{itemize}

\mediumspace

\textbf{میراث:} پیش‌درآمد تکنوکراسی و برنامه‌ریزی. پیروانش بعدها به جهات مختلفی رفتند—برخی سرمایه‌دار بزرگ شدند!
\end{personbox}

\subsection{شارل فوریه}

\begin{personbox}[شارل فوریه (۱۷۷۲-۱۸۳۷)]

\textbf{ملیت:} فرانسوی

\textbf{ایده محوری:} فالانستر (کمون‌های هماهنگ)

\mediumspace

\person{فوریه} از عجیب‌ترین و خلاق‌ترین متفکران سوسیالیست بود. او نظامی کامل برای جامعه آینده طراحی کرد:

\begin{itemize}[nosep]
    \item \textbf{فالانستر:} جوامع ۱۶۰۰ نفری خودکفا
    \item \textbf{کار جذاب:} کار باید لذت‌بخش باشد، نه عذاب
    \item \textbf{آزادی جنسی:} نقد خانواده سنتی و سرکوب جنسی
    \item \textbf{هماهنگی احساسات:} طراحی اجتماعی بر اساس روان‌شناسی
\end{itemize}

\mediumspace

\textbf{نقل‌قول:} \textit{«تمدن، دور باطل فقر است. کارگر فقیر است چون کار نمی‌کند؛ کار نمی‌کند چون فقیر است.»}

\mediumspace

\textbf{میراث:} چندین «فالانکس» در آمریکا تأسیس شد. فمینیست‌ها و طرفداران «کار لذت‌بخش» از او الهام گرفته‌اند.
\end{personbox}

\subsection{رابرت اوون}

\begin{personbox}[رابرت اوون (۱۷۷۱-۱۸۵۸)]

\textbf{ملیت:} ولزی-بریتانیایی

\textbf{ایده محوری:} تعاونی‌گری و اصلاح اجتماعی

\mediumspace

\person{اوون} برخلاف سن‌سیمون و فوریه، یک نظریه‌پرداز محض نبود—او صاحب کارخانه بود و ایده‌هایش را عملی کرد:

\begin{itemize}[nosep]
    \item در کارخانه نیو لانارک: کاهش ساعت کار، مهدکودک، آموزش کارگران
    \item تأسیس «نیو هارمونی» در آمریکا (شکست خورد)
    \item بنیان‌گذاری جنبش تعاونی بریتانیا
    \item تلاش برای ایجاد «بازار مبادله کار»
\end{itemize}

\mediumspace

\textbf{اعتقاد اصلی:} «محیط» شخصیت انسان را می‌سازد. با اصلاح محیط، می‌توان انسان بهتری ساخت.

\mediumspace

\textbf{میراث:} جنبش تعاونی جهانی، سوسیالیسم اصلاح‌طلب
\end{personbox}

\begin{historybox}[تعاونی‌های راچدیل (۱۸۴۴)]
در شهر راچدیل انگلستان، ۲۸ کارگر بافنده، نخستین تعاونی مصرف مدرن را بنیان نهادند. \textbf{اصول راچدیل}—مانند یک عضو یک رأی، سود به نسبت خرید، عضویت آزاد—هنوز مبنای تعاونی‌های جهان است. امروز بیش از ۱ میلیارد نفر عضو تعاونی‌ها هستند.
\end{historybox}

\subsection{نقد مارکس به سوسیالیسم تخیلی}

\begin{debatebox}[چرا مارکس آن‌ها را «تخیلی» نامید؟]
مارکس و انگلس در \book{مانیفست کمونیست} و جاهای دیگر، به این متفکران احترام گذاشتند اما آن‌ها را نقد کردند:

\textbf{نقدها:}
\begin{itemize}[nosep]
    \item آن‌ها جامعه آینده را «طراحی» می‌کردند، به جای تحلیل تضادهای جامعه موجود
    \item به طبقه کارگر به چشم «بیچاره‌های نیازمند کمک» نگاه می‌کردند، نه عامل تاریخی
    \item از راه‌حل‌های صلح‌آمیز و تدریجی دفاع می‌کردند، در حالی که تضاد طبقاتی راه‌حل صلح‌آمیز ندارد
    \item بر «ایده‌های خوب» تکیه می‌کردند، نه بر شرایط مادی
\end{itemize}

\textbf{اما امروز بسیاری می‌گویند:}
\begin{itemize}[nosep]
    \item شاید «تخیلی» بودن عیب نباشد—نیاز به آرمان داریم
    \item تجربه اوون نشان داد اصلاح تدریجی ممکن است
    \item توجه فوریه به روان‌شناسی و لذت، مهم بود
\end{itemize}
\end{debatebox}

% ══════════════════════════════════════════════════════════════════════════════
\section{سوسیالیسم علمی: مارکس و انگلس}
% ══════════════════════════════════════════════════════════════════════════════

\person{کارل مارکس} و \person{فردریش انگلس} ادعا کردند سوسیالیسم را از «آرمان‌شهر» به «علم» تبدیل کرده‌اند. این ادعا، چه درست باشد چه نه، تأثیری عظیم بر تاریخ گذاشت.

\subsection{مبانی مارکسیسم}

\begin{defbox}[تعریف: ماتریالیسم تاریخی (خلاصه)]
نظریه مارکس درباره تاریخ:
\begin{enumerate}[nosep]
    \item \textbf{زیربنا/روبنا:} شیوه تولید اقتصادی (زیربنا)، فرهنگ، حقوق، دین و سیاست (روبنا) را شکل می‌دهد
    \item \textbf{تضاد طبقاتی:} تاریخ، تاریخ مبارزه طبقات است
    \item \textbf{شیوه‌های تولید:} برده‌داری ← فئودالیسم ← سرمایه‌داری ← سوسیالیسم ← کمونیسم
    \item \textbf{تضاد درونی:} هر شیوه تولید تضادهایی دارد که آن را نابود می‌کند
    \item \textbf{انقلاب:} گذار از یک شیوه به دیگری معمولاً انقلابی است
\end{enumerate}
\end{defbox}

\begin{figure}[htbp]
\centering
\begin{tikzpicture}[scale=0.85]

% عنوان
\node[font=\large\bfseries, text=DarkGray] at (0, 7.5) {\fa{ماتریالیسم تاریخی: سیر تکامل جوامع}};

% مراحل
% برده‌داری
\fill[MediumGray, rounded corners=5pt] (-6.5, 4.5) rectangle (-4.5, 6);
\node[text=white, font=\bfseries] at (-5.5, 5.5) {\fa{برده‌داری}};
\node[text=VeryLightGray, font=\scriptsize] at (-5.5, 5) {\fa{برده‌دار/برده}};

% فلش
\draw[thickarrow, color=DarkGray] (-4.3, 5.25) -- (-3.2, 5.25);

% فئودالیسم
\fill[RightBlueLight, rounded corners=5pt] (-3, 4.5) rectangle (-1, 6);
\node[text=white, font=\bfseries] at (-2, 5.5) {\fa{فئودالیسم}};
\node[text=VeryLightGray, font=\scriptsize] at (-2, 5) {\fa{ارباب/رعیت}};

% فلش
\draw[thickarrow, color=DarkGray] (-0.8, 5.25) -- (0.3, 5.25);

% سرمایه‌داری
\fill[RightBlue, rounded corners=5pt] (0.5, 4.5) rectangle (2.5, 6);
\node[text=white, font=\bfseries] at (1.5, 5.5) {\fa{سرمایه‌داری}};
\node[text=VeryLightGray, font=\scriptsize] at (1.5, 5) {\fa{سرمایه‌دار/کارگر}};

% فلش انقلاب
\draw[thickarrow, color=LeftRed, line width=3pt] (2.7, 5.25) -- (3.8, 5.25);
\node[text=LeftRed, font=\tiny\bfseries] at (3.25, 5.7) {\fa{انقلاب!}};

% سوسیالیسم
\fill[LeftRedLight, rounded corners=5pt] (4, 4.5) rectangle (6, 6);
\node[text=white, font=\bfseries] at (5, 5.5) {\fa{سوسیالیسم}};
\node[text=VeryLightGray, font=\scriptsize] at (5, 5) {\fa{دیکتاتوری پرولتاریا}};

% فلش
\draw[thickarrow, color=LeftRed] (5, 4.3) -- (5, 3.2);

% کمونیسم
\fill[LeftRedDark, rounded corners=8pt] (3.5, 1.5) rectangle (6.5, 3);
\node[text=white, font=\bfseries\large] at (5, 2.5) {\fa{کمونیسم}};
\node[text=VeryLightGray, font=\scriptsize] at (5, 1.9) {\fa{جامعه بی‌طبقه}};

% توضیحات پایین
\node[text=MediumGray, font=\footnotesize, text width=10cm, align=center] at (0, 0.5) {
    \fa{در هر مرحله، طبقه‌ای علیه طبقه حاکم می‌شورد و نظم جدیدی می‌سازد.}\\
    \fa{سرمایه‌داری آخرین جامعه طبقاتی است؛ پرولتاریا آخرین طبقه انقلابی.}
};

\end{tikzpicture}
\caption{شیوه‌های تولید در نظریه مارکس}
\label{fig:historical-materialism}
\end{figure}

\subsection{نقد اقتصاد سیاسی}

مارکس در \book{سرمایه} (جلد اول: ۱۸۶۷) تحلیلی عمیق از سرمایه‌داری ارائه داد:

\begin{defbox}[مفاهیم کلیدی سرمایه]
\begin{itemize}[nosep]
    \item \textbf{کالا:} هر چیز تولیدشده برای مبادله (نه مصرف مستقیم)
    \item \textbf{ارزش مصرفی و ارزش مبادله‌ای:} مفید بودن vs قیمت
    \item \textbf{نیروی کار به مثابه کالا:} کارگر نیروی کارش را می‌فروشد
    \item \textbf{ارزش اضافی:} تفاوت بین ارزش تولیدشده و دستمزد
    \item \textbf{سرمایه:} پول که برای کسب سود بیشتر به کار می‌رود (P-C-P')
    \item \textbf{انباشت:} سرمایه باید مدام رشد کند، وگرنه نابود می‌شود
\end{itemize}
\end{defbox}

\begin{quotebox}[کارل مارکس، سرمایه]
\textit{«سرمایه، کار مُرده است که مانند خون‌آشام، فقط با مکیدن کار زنده زنده می‌ماند، و هرچه بیشتر بمکد، زنده‌تر است.»}
\end{quotebox}

\subsection{انقلاب و دیکتاتوری پرولتاریا}

\begin{defbox}[تعریف: دیکتاتوری پرولتاریا]
در نظریه مارکس:
\begin{itemize}[nosep]
    \item پس از انقلاب، طبقه کارگر قدرت را به دست می‌گیرد
    \item این دوره «گذار» است—دولت کارگری مقاومت بورژوازی را درهم می‌شکند
    \item با از بین رفتن طبقات، دولت نیز «زوال» می‌یابد
    \item نقطه پایان: جامعه بی‌طبقه و بی‌دولت (کمونیسم)
\end{itemize}
\textbf{توجه:} مارکس «دیکتاتوری» را به معنای حکومت یک طبقه به کار برد، نه لزوماً استبداد. اما این واژه بعدها توجیه‌گر استبداد شد.
\end{defbox}

% ══════════════════════════════════════════════════════════════════════════════
\section{شاخه‌های پس از مارکس}
% ══════════════════════════════════════════════════════════════════════════════

پس از مرگ مارکس (۱۸۸۳)، پیروانش در تفسیر آموزه‌هایش اختلاف پیدا کردند. این اختلافات به انشعابات بزرگی انجامید.

\subsection{انترناسیونال اول و دوم}

\begin{historybox}[انترناسیونال‌ها]
\textbf{انترناسیونال اول (۱۸۶۴-۱۸۷۶):}
\begin{itemize}[nosep]
    \item تأسیس با حضور مارکس در لندن
    \item تلاش برای اتحاد جنبش‌های کارگری
    \item جدال مارکس و باکونین (آنارشیست) → انشعاب
\end{itemize}

\textbf{انترناسیونال دوم (۱۸۸۹-۱۹۱۴):}
\begin{itemize}[nosep]
    \item احزاب سوسیالیست اروپا
    \item روز کارگر (اول ماه مه) از اینجا آمد
    \item رشد چشمگیر: SPD آلمان بزرگ‌ترین حزب شد
    \item فروپاشی با جنگ جهانی اول—احزاب از دولت‌های خود حمایت کردند!
\end{itemize}
\end{historybox}

\begin{figure}[htbp]
\centering
\begin{tikzpicture}[
    grow=down,
    level 1/.style={sibling distance=5.5cm, level distance=2cm},
    level 2/.style={sibling distance=2.8cm, level distance=2cm},
    edge from parent/.style={draw, line width=1.5pt, color=MediumGray},
    every node/.style={font=\small}
]

% ریشه
\node[rectangle, rounded corners=5pt, fill=LeftRedDark, text=white, font=\bfseries, minimum width=3cm, minimum height=0.8cm] {\fa{مارکسیسم}}
    child {
        node[rectangle, rounded corners=5pt, fill=LeftRed, text=white, minimum width=2.5cm, minimum height=0.7cm] {\fa{انقلابی}}
        child {
            node[rectangle, rounded corners=5pt, fill=LeftRedDark, text=white, minimum width=2cm, minimum height=0.6cm] {\fa{لنینیسم}}
        }
        child {
            node[rectangle, rounded corners=5pt, fill=OrangeAccent, text=white, minimum width=2cm, minimum height=0.6cm] {\fa{چپ کمونیست}}
        }
    }
    child {
        node[rectangle, rounded corners=5pt, fill=LeftRedLight, text=white, minimum width=2.5cm, minimum height=0.7cm] {\fa{اصلاح‌طلب}}
        child {
            node[rectangle, rounded corners=5pt, fill=CenterGreen, text=white, minimum width=2cm, minimum height=0.6cm] {\fa{سوسیال‌دموکراسی}}
        }
        child {
            node[rectangle, rounded corners=5pt, fill=CenterGreenLight, text=DarkGray, minimum width=2cm, minimum height=0.6cm] {\fa{فابیانیسم}}
        }
    };

% عنوان
\node[font=\large\bfseries, text=DarkGray] at (0, 1.5) {\fa{انشعاب بزرگ در مارکسیسم}};

% توضیح
\node[text=MediumGray, font=\footnotesize] at (0, -5) {\fa{پس از ۱۹۱۴ و به‌ویژه ۱۹۱۷، سوسیالیسم به دو شاخه اصلی تقسیم شد.}};

\end{tikzpicture}
\caption{شاخه‌های اصلی سوسیالیسم پس از مارکس}
\label{fig:socialism-branches}
\end{figure}

\subsection{انشعاب بزرگ: انقلابی در برابر اصلاح‌طلب}

مهم‌ترین انشعاب در تاریخ سوسیالیسم، جدال میان \keyword{انقلابیون} و \keyword{اصلاح‌طلبان} بود.

\begin{personbox}[ادوارد برنشتاین (۱۸۵۰-۱۹۳۲)]

\textbf{ملیت:} آلمانی

\textbf{اثر کلیدی:} \book{پیش‌شرط‌های سوسیالیسم} (۱۸۹۹)

\textbf{ایده محوری:} رویزیونیسم (بازنگری در مارکسیسم)

\mediumspace

\person{برنشتاین}، که با مارکس و انگلس آشنا بود، استدلال کرد:

\begin{itemize}[nosep]
    \item پیش‌بینی‌های مارکس درست درنیامده: طبقه متوسط محو نشده، بحران‌ها کمتر شده
    \item انقلاب خشونت‌آمیز نه ضروری است نه مطلوب
    \item سوسیالیسم باید از طریق \textbf{دموکراسی و اصلاحات} پیش برود
    \item «هدف نهایی هیچ است؛ جنبش همه چیز است»
\end{itemize}

\mediumspace

\textbf{واکنش:} حمله شدید از سوی کائوتسکی، لنین، رزا لوکزامبورگ. اما عملاً SPD همین راه را رفت.
\end{personbox}

\begin{personbox}[رزا لوکزامبورگ (۱۸۷۱-۱۹۱۹)]

\textbf{ملیت:} لهستانی-آلمانی

\textbf{آثار کلیدی:} \book{اصلاح یا انقلاب؟} (۱۹۰۰)، \book{انباشت سرمایه}

\textbf{ایده محوری:} مارکسیسم انقلابی و دموکراتیک

\mediumspace

\person{رزا لوکزامبورگ} هم با برنشتاین مخالفت کرد، هم بعدها با لنین:

\begin{itemize}[nosep]
    \item انقلاب ضروری است—اصلاحات در چارچوب سرمایه‌داری محدودند
    \item اما انقلاب باید \textbf{دموکراتیک} باشد، نه کودتای حزبی
    \item «آزادی همیشه آزادی دگراندیش است»
    \item نقد تمرکزگرایی لنین
\end{itemize}

\mediumspace

\textbf{سرنوشت:} در قیام اسپارتاکیست‌ها (۱۹۱۹) دستگیر و توسط نیروهای راست کشته شد.
\end{personbox}

\begin{comparebox}[مقایسه: برنشتاین، لنین، لوکزامبورگ]
\begin{table}[H]
\centering
\small
\begin{tabular}{r|c|c|c}
\toprule
\textbf{موضوع} & \textbf{برنشتاین} & \textbf{لنین} & \textbf{لوکزامبورگ} \\
\midrule
روش تغییر & اصلاحات تدریجی & انقلاب حزبی & انقلاب توده‌ای \\
نقش حزب & ابزار انتخاباتی & پیشاهنگ انقلاب & هماهنگ‌کننده \\
دموکراسی & هدف و ابزار & ابزار موقت & اصل غیرقابل چشم‌پوشی \\
دولت پس از انقلاب & دولت رفاه & دیکتاتوری حزب & دموکراسی شورایی \\
\bottomrule
\end{tabular}
\end{table}
\end{comparebox}

% ══════════════════════════════════════════════════════════════════════════════
\section{سوسیالیسم دموکراتیک}
% ══════════════════════════════════════════════════════════════════════════════

شاخه‌ای از سوسیالیسم که بر دموکراسی سیاسی تأکید دارد و انقلاب خشونت‌آمیز را رد می‌کند.

\subsection{فابیان‌ها در انگلستان}

\begin{historybox}[جامعه فابیان (تأسیس ۱۸۸۴)]
گروهی از روشنفکران بریتانیایی که استراتژی «تدریج‌گرایی» را برگزیدند:
\begin{itemize}[nosep]
    \item نام‌گذاری به یاد فابیوس ماکسیموس، ژنرال رومی که با صبر پیروز شد
    \item اعضای برجسته: جورج برنارد شاو، سیدنی و بئاتریس وب، اچ.جی. ولز
    \item تأثیر بر تأسیس حزب کارگر بریتانیا (۱۹۰۰)
    \item تأسیس مدرسه اقتصاد لندن (LSE)
\end{itemize}
\textbf{شعار:} «وقتی حمله کنی، بزن؛ وقتی نه، صبر کن.»
\end{historybox}

\subsection{SPD آلمان}

\begin{historybox}[حزب سوسیال‌دموکرات آلمان (SPD)]
بزرگ‌ترین و مهم‌ترین حزب سوسیالیست پیش از جنگ جهانی اول:
\begin{itemize}[nosep]
    \item تأسیس: ۱۸۷۵ (ادغام دو حزب)
    \item برنامه ارفورت (۱۸۹۱): مارکسیسم رسمی + اصلاحات عملی
    \item ۱۹۱۲: بزرگ‌ترین حزب رایشتاگ با ۳۴٪ آرا
    \item ۱۹۱۴: رأی به اعتبارات جنگی → شکاف با چپ رادیکال
    \item ۱۹۱۸-۱۹: انقلاب آلمان، جمهوری وایمار
\end{itemize}
\end{historybox}

% ══════════════════════════════════════════════════════════════════════════════
\section{سوسیالیسم گیلدی و سندیکالیسم}
% ══════════════════════════════════════════════════════════════════════════════

همه سوسیالیست‌ها دولت‌گرا نبودند. برخی بر سازمان‌دهی در محل کار تأکید داشتند.

\begin{defbox}[تعریف: سندیکالیسم]
جنبشی در اواخر قرن نوزدهم و اوایل بیستم که معتقد بود:
\begin{itemize}[nosep]
    \item \textbf{اتحادیه‌های کارگری} (سندیکاها) ابزار اصلی مبارزه و سازمان‌دهی جامعه آینده‌اند
    \item \textbf{اقدام مستقیم} (اعتصاب، خرابکاری) در برابر سیاست پارلمانی
    \item \textbf{اعتصاب عمومی} می‌تواند سرمایه‌داری را فلج و سرنگون کند
    \item دولت کارگری هم نباید باشد—کارگران خودشان کارخانه‌ها را اداره کنند
\end{itemize}
\textbf{نمونه:} CGT فرانسه، IWW آمریکا، CNT اسپانیا
\end{defbox}

\begin{defbox}[تعریف: سوسیالیسم گیلدی]
نسخه بریتانیایی و ملایم‌تر که ترکیب می‌کرد:
\begin{itemize}[nosep]
    \item خودگردانی کارگری در هر صنعت («گیلد»)
    \item دولت به عنوان نماینده مصرف‌کنندگان
    \item دموکراسی صنعتی + دموکراسی سیاسی
\end{itemize}
\textbf{چهره کلیدی:} جی.دی.اچ. کول
\end{defbox}

% ══════════════════════════════════════════════════════════════════════════════
\section{سوسیالیسم مسیحی و مذهبی}
% ══════════════════════════════════════════════════════════════════════════════

رابطه سوسیالیسم و مذهب پیچیده بوده است. مارکس دین را «افیون توده‌ها» خواند، اما همیشه گرایش‌های سوسیالیستی مذهبی هم وجود داشته‌اند.

\subsection{سوسیالیسم مسیحی}

\begin{defbox}[تعریف: سوسیالیسم مسیحی]
جنبشی که استدلال می‌کند:
\begin{itemize}[nosep]
    \item تعالیم مسیح (عدالت، کمک به فقرا، برابری) با سوسیالیسم سازگارند
    \item کلیسای اولیه نوعی اشتراک داشت
    \item سرمایه‌داری با اخلاق مسیحی ناسازگار است
\end{itemize}
\textbf{نمونه‌ها:} چارلز کینگزلی (انگلستان)، دوروتی دی (آمریکا)
\end{defbox}

\subsection{الهیات رهایی‌بخش}

\begin{historybox}[الهیات رهایی‌بخش (دهه ۱۹۶۰ به بعد)]
جنبشی در کلیسای کاتولیک آمریکای لاتین:
\begin{itemize}[nosep]
    \item گوستاوو گوتیرز: \book{الهیات رهایی‌بخش} (۱۹۷۱)
    \item «انتخاب ترجیحی فقرا»
    \item استفاده از تحلیل مارکسی (بدون الحاد)
    \item مقاومت در برابر دیکتاتوری‌های نظامی
    \item مخالفت واتیکان در دوره ژان پل دوم
\end{itemize}
\end{historybox}

\subsection{سوسیالیسم اسلامی؟}

\begin{debatebox}[آیا سوسیالیسم اسلامی ممکن است؟]
\textbf{موافقان:}
\begin{itemize}[nosep]
    \item زکات و انفاق نوعی توزیع ثروت است
    \item اسلام ربا (بهره) را حرام می‌داند—نقد سرمایه مالی
    \item «الناس شرکاء فی ثلاث: الماء و الکلاء و النار» (حدیث)
    \item نمونه‌ها: علی شریعتی، قذافی، چپ‌های مسلمان
\end{itemize}

\textbf{مخالفان:}
\begin{itemize}[nosep]
    \item اسلام مالکیت خصوصی را به رسمیت می‌شناسد
    \item سوسیالیسم معمولاً سکولار یا ضددینی بوده
    \item «سوسیالیسم اسلامی» اغلب اقتدارگرا بوده (مصر ناصر، لیبی)
\end{itemize}
\end{debatebox}

% ══════════════════════════════════════════════════════════════════════════════
\section{اختلافات کلیدی درون سوسیالیسم}
% ══════════════════════════════════════════════════════════════════════════════

همان‌طور که دیدیم، سوسیالیسم هرگز یکدست نبوده است. جدول زیر اختلافات اصلی را خلاصه می‌کند.

\begin{landscape}
\begin{table}[htbp]
\centering
\caption{اختلافات کلیدی درون سوسیالیسم}
\label{tab:socialism-debates}
\small
\begin{tabular}{r|c|c|c|c}
\toprule
\textbf{موضوع} & \textbf{یک قطب} & \textbf{قطب دیگر} & \textbf{نمونه قطب اول} & \textbf{نمونه قطب دوم} \\
\midrule
روش تغییر & انقلاب & اصلاح تدریجی & لنین، لوکزامبورگ & برنشتاین، فابیان‌ها \\
\midrule
نقش حزب & پیشاهنگ متمرکز & حزب توده‌ای & لنین & کائوتسکی \\
\midrule
دولت & ابزار رهایی & خطر (حتی دولت کارگری) & مارکسیست‌ها & آنارشیست‌ها \\
\midrule
بین‌الملل‌گرایی & کارگران همه کشورها متحد & ملی‌گرایی چپ & تروتسکی & استالین \\
\midrule
پایگاه طبقاتی & پرولتاریای صنعتی & دهقانان و خرده‌بورژوازی & مارکسیسم ارتدکس & مائو، پوپولیست‌ها \\
\midrule
دموکراسی & ابزار (تابع انقلاب) & اصل (غیرقابل چشم‌پوشی) & لنین (پس از ۱۹۱۷) & لوکزامبورگ \\
\midrule
بازار & باید محو شود & می‌تواند با سوسیالیسم باشد & برنامه‌ریزان & سوسیالیسم بازار \\
\midrule
مذهب & افیون توده‌ها & منبع عدالت‌خواهی & مارکسیسم ارتدکس & الهیات رهایی‌بخش \\
\bottomrule
\end{tabular}
\end{table}
\end{landscape}

% ══════════════════════════════════════════════════════════════════════════════
%                         خلاصه و پرسش‌ها
% ══════════════════════════════════════════════════════════════════════════════

\newpage

\begin{summarybox}

\textbf{۱. تعریف:} سوسیالیسم مجموعه‌ای متنوع است که در نقد سرمایه‌داری، تأکید بر برابری و همبستگی، و خواست مالکیت جمعی اشتراک دارند.

\textbf{۲. پیشگامان:} سن‌سیمون (تکنوکراسی)، فوریه (فالانستر)، اوون (تعاونی‌ها) ایده‌های اولیه را ارائه دادند. مارکس آن‌ها را «تخیلی» نامید.

\textbf{۳. مارکسیسم:} ماتریالیسم تاریخی، نقد اقتصاد سیاسی، ارزش اضافی، انقلاب پرولتری، دیکتاتوری پرولتاریا.

\textbf{۴. انشعاب بزرگ:} پس از مارکس، سوسیالیسم به دو شاخه اصلی تقسیم شد: انقلابی (لنینیسم) و اصلاح‌طلب (سوسیال‌دموکراسی).

\textbf{۵. سوسیالیسم دموکراتیک:} فابیان‌ها و SPD راه تدریجی و پارلمانی را برگزیدند.

\textbf{۶. گرایش‌های دیگر:} سندیکالیسم (اتحادیه‌محور)، سوسیالیسم گیلدی، سوسیالیسم مسیحی.

\textbf{۷. تنوع درونی:} سوسیالیست‌ها در مورد روش (انقلاب/اصلاح)، دولت، دموکراسی، بازار و مذهب اختلاف دارند.

\end{summarybox}

\vspace{20pt}

\begin{questionbox}

\begin{enumerate}[nosep, rightmargin=1cm]
    \item آیا نقد مارکس به سوسیالیست‌های «تخیلی» منصفانه بود؟ آیا «آرمان‌شهرگرایی» همیشه بد است؟

    \item برنشتاین گفت پیش‌بینی‌های مارکس درست درنیامده. با توجه به یک قرن بعدی، آیا حق داشت؟

    \item آیا می‌توان سوسیالیست بود و همزمان از دموکراسی لیبرال دفاع کرد؟ یا این ترکیب ناپایدار است؟

    \item رابطه سوسیالیسم و مذهب چگونه باید باشد؟ آیا «سوسیالیسم دینی» تناقض‌آمیز است؟
\end{enumerate}

\end{questionbox}

% ══════════════════════════════════════════════════════════════════════════════
%                         پایان فصل چهارم
% ══════════════════════════════════════════════════════════════════════════════